\documentclass[a4paper]{article}
% Español (México) con XeLaTeX: fontspec para fuentes Unicode, polyglossia
% para la división silábica y los nombres del documento en la variante mexicana.
\usepackage{fontspec}
\usepackage{polyglossia}
\setdefaultlanguage[variant=mexican]{spanish}
% Los nombres chinos van como «pinyin (original)»: xeCJK da a los caracteres
% Han su propia fuente; sin ella XeLaTeX los omite con un aviso.
\usepackage{xeCJK}
\setCJKmainfont{FandolSong-Regular.otf}
% polyglossia no traduce los nombres de listings.
\AtBeginDocument{\providecommand{\lstlistingname}{}\renewcommand{\lstlistingname}{Listado}%
  \providecommand{\lstlistlistingname}{}\renewcommand{\lstlistlistingname}{Índice de listados}}
\usepackage{amsmath, amssymb}
\usepackage{graphicx}
\usepackage[margin=2.5cm]{geometry}
\usepackage[most]{tcolorbox}
\usepackage{etoolbox}
\usepackage{listings}
\usepackage{hyperref}
\usepackage{booktabs}
\usepackage{subcaption}
\usepackage{float}

\newtcolorbox{knowledgebox}[1]{enhanced, colback=blue!5!white, colframe=blue!75!black, colbacktitle=blue!75!black, coltitle=white, fonttitle=\bfseries, title={#1}, attach boxed title to top left={yshift=-2mm, xshift=2mm}, boxrule=1pt, sharp corners}
\newtcolorbox{importantbox}[1]{enhanced, colback=yellow!10!white, colframe=yellow!80!black, colbacktitle=yellow!80!black, coltitle=black, fonttitle=\bfseries, title={#1}, sharp corners}
\newtcolorbox{warningbox}[1]{enhanced, colback=red!5!white, colframe=red!75!black, colbacktitle=red!75!black, coltitle=white, fonttitle=\bfseries, title={#1}, sharp corners}

\lstset{language=Python, basicstyle=\ttfamily\small, keywordstyle=\color{blue}, stringstyle=\color{red!60!black}, commentstyle=\color{green!60!black}, breaklines=true, frame=single, numbers=left, numberstyle=\tiny\color{gray}, captionpos=b, extendedchars=true}

\newcommand{\notetitle}{Salida estructurada y llamadas a herramientas: LangChain / OpenAI API}
\newcommand{\noteauthors}{Elaboradas a partir de materiales públicos del curso}
\newcommand{\notedate}{2025}
\newcommand{\videochannel}{Wudaokou Nashi (五道口纳什)}
\newcommand{\videopublishdate}{2025}
\newcommand{\videoduration}{aproximadamente 9 minutos}
\newcommand{\videourl}{https://www.bilibili.com/video/BV1uiUnBFED4}
\newcommand{\videocoverpath}{cover.jpg}
\newcommand{\repourl}{https://github.com/hqhq1025/ai-course-notes}


% Footer with repo info
\usepackage{fancyhdr}
\pagestyle{fancy}
\fancyhf{}
\fancyfoot[C]{\thepage}
\fancyfoot[R]{\footnotesize\color{gray}\href{\repourl}{AI Course Notes}}
\renewcommand{\headrulewidth}{0pt}
\renewcommand{\footrulewidth}{0pt}

\begin{document}
\begin{titlepage}
\centering
{\Large Notas del curso\par}\vspace{1.2cm}
{\huge\bfseries \notetitle\par}\vspace{0.8cm}
{\large \noteauthors\par}\vspace{0.3cm}
{\large \notedate\par}\vspace{1.2cm}
\ifdefempty{\videocoverpath}{}{\includegraphics[width=0.82\textwidth,height=0.45\textheight,keepaspectratio]{\videocoverpath}\par}
\vfill
\begin{tcolorbox}[width=0.9\textwidth, colback=black!2!white, colframe=black!60, sharp corners]
\textbf{Autor o canal del video}: \videochannel\par\textbf{Fecha de publicación}: \videopublishdate\par\textbf{Duración del video}: \videoduration\par
\ifdefempty{\videourl}{}{\textbf{Enlace del video}: \href{\videourl}{\nolinkurl{\videourl}}\par}\par
\textbf{Página del proyecto}: \href{\repourl}{\nolinkurl{\repourl}}\par
\end{tcolorbox}
\end{titlepage}
\tableofcontents\newpage

\section{La naturaleza de la salida estructurada y las llamadas a herramientas}

\subsection{Por qué se necesita la salida estructurada}

Al construir Agent Workflows complejos, los módulos necesitan un \textbf{acuerdo de interfaz} explícito: la salida de una etapa es la entrada de otra. La salida estructurada resuelve tres problemas centrales:

\begin{enumerate}
    \item No se necesita acordar el Output Format a nivel de Prompt
    \item No se necesita analizar manualmente JSON complejo
    \item La información de Function / Tools tampoco necesita colocarse explícitamente en el Prompt
\end{enumerate}

\begin{importantbox}{Salida estructurada = SDD en versión débil}
La salida estructurada implementa, en cierto sentido, una \textbf{versión débil del Specification Driven Development}: al definir el formato de entrada y salida entre los Módulos, simplifica enormemente el trabajo de análisis y adaptación. La razón subyacente es que los modelos de lenguaje grandes actuales ya están muy bien entrenados en la salida estructurada.
\end{importantbox}

\subsection{Resumen de la sección}

La salida estructurada y las llamadas a herramientas son la infraestructura del desarrollo de Agentes: mediante un acuerdo de formato a nivel de API, se omite una gran cantidad de trabajo manual de diseño de Prompt y análisis de la salida.

\section{Implementación a nivel de la API de OpenAI}

\subsection{Structured Output: el campo response\_format}

Se pasa el JSON Schema del modelo de Pydantic a través del campo \texttt{response\_format}:

\begin{itemize}
    \item Se define un modelo de Pydantic (por ejemplo, con los campos title y year)
    \item Se serializa a JSON Schema mediante \texttt{model.json\_schema()}
    \item Se coloca en el campo \texttt{response\_format} de la solicitud HTTP
    \item La salida del modelo se ajusta automáticamente a ese Schema y puede deserializarse directamente en un objeto
\end{itemize}

\subsection{Function Calling: el campo tools}

Se inyecta la información de las herramientas mediante el campo \texttt{tools}:

\begin{itemize}
    \item Se registra la función de Python con el decorador \texttt{@tool} de LangChain
    \item Se convierte al formato que exige la API de OpenAI mediante \texttt{convert\_to\_openai\_tool()}
    \item Se coloca en el campo \texttt{tools} de la solicitud HTTP
    \item La respuesta del modelo incluye el contenido instanciado de \texttt{function\_call}
\end{itemize}

\begin{knowledgebox}{Análisis por captura de tráfico HTTP}
Mediante la captura de tráfico HTTP se puede ver con claridad que:
\begin{itemize}
    \item \textbf{Structured Output} pasa por el campo \texttt{response\_format}
    \item \textbf{Function Calling / Bound Tools} pasa por el campo \texttt{tools}
\end{itemize}
Ambos son campos distintos a nivel de la solicitud HTTP, pero su propósito es el mismo: fijar el formato de entrada y salida.
\end{knowledgebox}

\subsection{Resumen de la sección}

A nivel de la solicitud HTTP de la API de OpenAI, la salida estructurada se implementa mediante response\_format y la llamada a herramientas mediante tools; entender el mecanismo subyacente ayuda a depurar y optimizar.

\section{Encapsulamiento a nivel de LangChain}

\subsection{with\_structured\_output}

El método \texttt{with\_structured\_output} de LangChain simplifica aún más el uso:

\begin{itemize}
    \item Se pasa directamente un Pydantic Model
    \item Internamente se convierte de manera automática a JSON Schema y se coloca en response\_format
    \item El resultado devuelto ya es un objeto instanciado (un paso más de deserialización automática respecto a la API de OpenAI)
\end{itemize}

\begin{warningbox}{Diferencia entre las dos rutas}
\begin{itemize}
    \item \texttt{with\_structured\_output} $\rightarrow$ pasa por \texttt{response\_format} (el JSON Schema se alinea con la semántica de los campos)
    \item \texttt{bind\_tools} $\rightarrow$ pasa por \texttt{tools} (semántica de Function Calling)
\end{itemize}
Un Pydantic Model es, en esencia, también una especie de Function (instanciar la clase $\approx$ llamar al constructor), de modo que bind\_tools de LangChain también puede aceptar un Pydantic Model.
\end{warningbox}

\subsection{Resumen de la sección}

LangChain agrega un encapsulamiento adicional sobre la API de OpenAI: serialización automática, deserialización automática, con una experiencia de uso más simple.

\section{Técnica práctica: implementar CoT mediante salida estructurada}

\begin{importantbox}{Implementar Chain of Thought mediante salida estructurada}
El proceso de razonamiento se puede encapsular como un objeto de salida estructurada de Pydantic, por ejemplo:
\begin{itemize}
    \item \texttt{steps}: {[}\texttt{Step(explanation=..., output=...)}{]}
    \item Cada Step contiene una explicación y una salida
\end{itemize}
Al convenir la implementación del efecto CoT de esta manera, se puede lograr de forma natural una mejora considerable en el desempeño del modelo. Esta es una técnica muy práctica para el uso de modelos de lenguaje.
\end{importantbox}

\section{Síntesis y ampliación}

\begin{enumerate}
    \item \textbf{La salida estructurada y las llamadas a herramientas son la infraestructura del desarrollo de Agentes}: implementan, a nivel de la API, el acuerdo de interfaz entre módulos.
    \item \textbf{Entender la solicitud HTTP subyacente} ayuda a depurar: response\_format vs. tools son dos campos distintos.
    \item \textbf{LangChain ofrece un encapsulamiento de más alto nivel}: serialización/deserialización automáticas, pero por debajo sigue siendo la API de OpenAI.
    \item \textbf{La salida estructurada permite implementar CoT}: estructurar el proceso de razonamiento es una técnica práctica para mejorar el desempeño.
    \item \textbf{La capacidad de salida estructurada de los LLM modernos ya es muy sólida}: puede prescindirse de una gran cantidad de convenciones manuales de formato en el prompt y de análisis manual de la salida.
\end{enumerate}

\end{document}
